\documentclass[../main.tex]{subfiles}
\begin{document}

\section{Discussion}
\label{sec:discussion}

\subsection{From heuristic to criterion}

The practical value of Corollary~\ref{cor:critical} is that it converts a design
intuition into a quantity an operator can estimate. The inputs are the mean
neighbourhood size $D$, the bot prevalence $\pi_B$, and the feature signal-to-noise
ratio $\|\Delta\|/\sigma$---all measurable on a deployed system---and the output is
a threshold on camouflage beyond which the graph is a liability. This reframes a
question usually settled by ablation (``does the graph help on this benchmark?'')
as one that can be answered per-cohort: for hub accounts with large $D$ the graph
remains valuable at camouflage levels that make it actively harmful for sparse
accounts. A detector that applies one aggregation policy to both is leaving
accuracy on the table at one end and incurring damage at the other.

The analysis also isolates an under-appreciated coupling. Because
$\eta = c\rho$, bot prevalence enters the critical ratio directly: a platform whose
bot population doubles does not merely face more bots, it faces a detector whose
tolerance to each bot's camouflage has fallen. Class imbalance is normally treated
as a training-time nuisance to be fixed by resampling; \eqref{eq:critical} shows it
is also a structural robustness parameter, and resampling does not touch it.

\subsection{What the theory does and does not license}

We are deliberate about the boundaries of these claims.

The results characterise \emph{one round of mean aggregation} in a stylised
generative model. They explain why homophilic message passing fails under
camouflage and why the three mitigations help, and they give the correct
qualitative orderings---more pruning is better, larger neighbourhoods are safer,
higher prevalence is more dangerous. They do not predict the accuracy of a trained
multi-layer network on a real bot graph, and no number in
Section~\ref{sec:validation} should be read as such. Our simulations verify
algebra; they are not empirical evidence about Twitter or any other platform.

Three specific gaps deserve naming. First, the theory is single-layer; depth
introduces over-smoothing effects~\cite{li2018deeper,oono2020graph} that compound
with camouflage in ways this analysis does not capture. Second, we do not model
learning: the SGTR attenuation $\bar e_c$ is treated as given, whereas in practice
it is the output of a trained scorer whose quality depends on the very separability
the pruning is meant to restore. Closing that loop---bounding $\bar e_c$ as a
function of $\|\Delta\|$, $d$, and sample size---is the natural next step and would
turn Theorem~\ref{thm:sgtr} from a conditional into an unconditional guarantee.
Third, the disentanglement analysis is the weakest link, and we have said why: the
orthogonality penalty controls dependence only in the Gaussian case
(Proposition~\ref{prop:mi}), and identifiability rests entirely on label
supervision.

\subsection{Limitations of the generative model}
\label{sec:limitations}

The camouflage-CSBM makes four simplifications, each of which trades realism for
tractability.

\emph{Isotropic, equal-covariance features.} Real account features are neither
isotropic nor identically dispersed across classes. Anisotropy would replace
$\|\Delta\|/\sigma$ with a Mahalanobis distance throughout; the structure of the
results survives, but $\rho^\star$ becomes direction-dependent.

\emph{Fixed neighbourhood size.} Social graphs are heavy-tailed. Since $\rho^\star$
is concave in $D$, a mixture of degrees is \emph{not} well summarised by its mean:
applying~\eqref{eq:critical} at the average degree overstates robustness for the
low-degree majority. Per-degree stratification, which the theory directly supports,
is the appropriate remedy.

\emph{Independent neighbour sampling.} We ignore triadic closure and community
structure, both prominent in real follow graphs. Correlated neighbourhoods reduce
the effective sample size in the averaging step, which inflates the variance
in~\eqref{eq:variances} and lowers $\rho^\star$; our thresholds are therefore
optimistic.

\emph{A single homogeneous relation.} Production detectors operate on
heterogeneous graphs with several edge types. The analysis extends per-relation,
and the interesting question---how an adversary should allocate a fixed camouflage
budget across relations---follows naturally from~\eqref{eq:critical} but is beyond
the present scope.

\subsection{An adversarial reading}

Finally, the theory can be read from the attacker's side, which we think is the
more useful test of it. Corollary~\ref{cor:critical} tells an adversary exactly how
much camouflage is needed to neutralise a homophilic detector, and
Theorem~\ref{thm:sgtr} tells them how much more is needed against a pruning
defence. That $\rho^\star_{\mathrm{SGTR}} \to 1$ only as $\bar e_c \to 1$ means no
imperfect pruning rule is a permanent fix; the defence buys operating range, and an
adversary willing to spend edges can always exceed it. Treating $\rho$ as a
strategic variable rather than a nuisance parameter---and analysing the resulting
game between detector and adversary---strikes us as the most promising direction
this line of work opens.

\section{Conclusion}
\label{sec:conclusion}

We have given an analytical account of relation camouflage in GNN-based bot
detection. Modelling camouflage as an adversarial parameter in a contextual
stochastic block model, we proved that mean aggregation deflates the class signal
by exactly $1-(1+c)\rho$, derived the resulting Bayes error, and obtained a
critical camouflage ratio $\rho^\star = (1-D^{-1/2})/(1+c)$ beyond which a
homophilic GNN is provably worse than ignoring the graph---a result that recovers
the known $1/\sqrt{D}$ separability gain of graph convolution at $\rho = 0$. We
then showed that spectral edge pruning separates camouflage edges by exactly
$\|\Delta\|^2$ in expectation and raises the tolerated camouflage to
$\rho^\star/(\rho^\star + (1-\bar e_c)(1-\rho^\star))$; that tri-channel
aggregation is strictly more expressive than any single homophilic aggregator; and
that the contrastive prototype objective maximises a variational lower bound on
label information while its orthogonality penalty guarantees independence only
under a joint-Gaussian assumption. All closed forms were verified numerically, with
the predicted critical ratio falling inside the measured crossover bracket.

The contribution is not a new detector but a criterion. Practitioners currently
decide whether to trust a graph-based detector by running an ablation on a
benchmark; the results here let them instead estimate three measurable quantities
and compute the answer for their own deployment, per account cohort. Extending the
analysis to multiple layers, to heterogeneous relations, and to a learned rather
than assumed pruning quality would make that criterion sharper still.

\end{document}
